\documentclass[10pt,a4paper]{amsart}
\usepackage{amsmath,amssymb,amsthm,graphicx}
% Retained preamble for the Stacks Project chapter snapshot in this directory.
% The chapter's own preamble.tex is a copy of the shared Stacks Project preamble; the block
% below is sliced MECHANICALLY, line for line, from it. Nothing was added, removed or
% redefined; the wrapper class, the xr-hyper/hyperref wiring and the externaldocument list
% are omitted because the assembler supplies the document class.
%
% lines taken from preamble.tex: 3, 5, 9, 11, 12, 156, 157, 158, ... (51 lines total)
%
% package coverage restored from the chapter's own preamble (2026-10-13): xy, xy-2cell

\usepackage[all]{xy}
\xyoption{2cell}
\UseAllTwocells

\newenvironment{reference}{\comment}{\endcomment}
\newenvironment{slogan}{\comment}{\endcomment}
\newenvironment{history}{\comment}{\endcomment}
\theoremstyle{plain}
\newtheorem{theorem}[subsection]{Theorem}
\newtheorem{proposition}[subsection]{Proposition}
\newtheorem{lemma}[subsection]{Lemma}
\theoremstyle{definition}
\newtheorem{definition}[subsection]{Definition}
\newtheorem{example}[subsection]{Example}
\newtheorem{exercise}[subsection]{Exercise}
\newtheorem{situation}[subsection]{Situation}
\theoremstyle{remark}
\newtheorem{remark}[subsection]{Remark}
\newtheorem{remarks}[subsection]{Remarks}
\numberwithin{equation}{subsection}
\def\lim{\mathop{\mathrm{lim}}\nolimits}
\def\colim{\mathop{\mathrm{colim}}\nolimits}
\def\Spec{\mathop{\mathrm{Spec}}}
\def\Hom{\mathop{\mathrm{Hom}}\nolimits}
\def\Ext{\mathop{\mathrm{Ext}}\nolimits}
\def\SheafHom{\mathop{\mathcal{H}\!\mathit{om}}\nolimits}
\def\SheafExt{\mathop{\mathcal{E}\!\mathit{xt}}\nolimits}
\def\Sch{\mathit{Sch}}
\def\Mor{\mathop{\mathrm{Mor}}\nolimits}
\def\Ob{\mathop{\mathrm{Ob}}\nolimits}
\def\Sh{\mathop{\mathit{Sh}}\nolimits}
\def\NL{\mathop{N\!L}\nolimits}
\def\CH{\mathop{\mathrm{CH}}\nolimits}
\def\proetale{{pro\text{-}\acute{e}tale}}
\def\etale{{\acute{e}tale}}
\def\QCoh{\mathit{QCoh}}
\def\Ker{\mathop{\mathrm{Ker}}}
\def\Im{\mathop{\mathrm{Im}}}
\def\Coker{\mathop{\mathrm{Coker}}}
\def\Coim{\mathop{\mathrm{Coim}}}
\DeclareMathSymbol{\boxtimes}{\mathbin}{AMSa}{"02}
\def\QCohstack{\mathcal{QC}\!\mathit{oh}}
\def\Cohstack{\mathcal{C}\!\mathit{oh}}
\def\Spacesstack{\mathcal{S}\!\mathit{paces}}
\def\Quotfunctor{\mathrm{Quot}}
\def\Hilbfunctor{\mathrm{Hilb}}
\def\Curvesstack{\mathcal{C}\!\mathit{urves}}
\def\Polarizedstack{\mathcal{P}\!\mathit{olarized}}
\def\Complexesstack{\mathcal{C}\!\mathit{omplexes}}
\def\Pic{\mathop{\mathrm{Pic}}\nolimits}
\def\Picardstack{\mathcal{P}\!\mathit{ic}}
\def\Picardfunctor{\mathrm{Pic}}
\def\Deformationcategory{\mathcal{D}\!\mathit{ef}}
\usepackage{cleveref}
\crefname{originallemma}{original lemma}{original lemmas}
\Crefname{originallemma}{Original lemma}{Original lemmas}
\makeatletter
\newlabel{lemma-coherent-functorial}{{0DN9}{source}{}{sourcepointer-1}{}}
\newlabel{lemma-coherent-functorial@cref}{{[originallemma][0][]0DN9}{[1][0][]source}{}{}{}}
\makeatother
\crefname{originallemma}{original lemma}{original lemmas}
\Crefname{originallemma}{Original lemma}{Original lemmas}
\makeatletter
\newlabel{spaces-flat-lemma-F-zero-closed-proper}{{083M}{source}{}{sourcepointer-2}{}}
\newlabel{spaces-flat-lemma-F-zero-closed-proper@cref}{{[originallemma][0][]083M}{[1][0][]source}{}{}{}}
\makeatother
\crefname{originallemma}{original lemma}{original lemmas}
\Crefname{originallemma}{Original lemma}{Original lemmas}
\makeatletter
\newlabel{spaces-morphisms-lemma-i-star-equivalence}{{04CJ}{source}{}{sourcepointer-3}{}}
\newlabel{spaces-morphisms-lemma-i-star-equivalence@cref}{{[originallemma][0][]04CJ}{[1][0][]source}{}{}{}}
\makeatother
\title{Closed immersions of bases induce closed immersions of Coh stacks}
\author{The Stacks Project Authors}
\date{stacks-project:moduli.tex:local-snapshot-20261012; GNU FDL 1.2 or later}
\begin{document}
\maketitle
\noindent\textit{Editorial note.} Statement span is lines 511-517 of sources/stacks-moduli/pinned-original-source.tex; proof span is lines 519-552, the complete author proof (32 body lines).\par\medskip
\noindent\textit{Editorial note.} Navigation stubs: 3 cited-but-not-included labels are supplied as non-load-bearing pointers with their pinned tags where the snapshot tag index has them.\par\medskip
\noindent\textit{Editorial note.} extra\_packages is empty on purpose: the retained preamble already loads xy with options and a second declaration clashes in the assembler wrapper.\par\medskip
\noindent\textit{Editorial note.} Text is retained verbatim, including the author's own typos and placeholders.\par\medskip
\noindent\textit{Editorial note.} Queue index 12 of queue\_moduli.json (pb=1414, stmt 511-517, proof 519-552).\par\medskip
\noindent\textit{Editorial note.} stacks-moduli wave-1: fresh adjudication accept (2026-10-14); pool per Chief order 08:5x (dpa superseded to lane-s9); moduli.tex 53 labels/48 with proof judged 25 accept / 23 reject / 5 pb=0 env-dup rows excluded; spans cite-free and reference/slogan/history-free so no patchenv needed.\par\medskip
\noindent\textit{Allowed standard prerequisites (editorial).}\begin{itemize}
\item Explicit prerequisite boundary. The cited inputs are prerequisites; the argument in the retained span is not replaced by any cited result.
\end{itemize}\medskip
\ifdefined\hypertarget\hypertarget{sourcepointer-1}{}\fi
\noindent\textit{Source reference.} Original lemma 0DN9: sources/stacks-moduli/pinned-original-source.tex, source lines 425--430. Cited by the retained statement or proof but not part of the collected spans; supplied only as a navigation pointer. The argument does not depend on its content.\par\medskip
\ifdefined\hypertarget\hypertarget{sourcepointer-2}{}\fi
\noindent\textit{Source reference.} Original lemma 083M: sources/stacks-moduli/pinned-original-source.tex, source lines 552--552. Cited by the retained statement or proof but not part of the collected spans; supplied only as a navigation pointer. The argument does not depend on its content.\par\medskip
\ifdefined\hypertarget\hypertarget{sourcepointer-3}{}\fi
\noindent\textit{Source reference.} Original lemma 04CJ: sources/stacks-moduli/pinned-original-source.tex, source lines 552--552. Cited by the retained statement or proof but not part of the collected spans; supplied only as a navigation pointer. The argument does not depend on its content.\par\medskip
\begin{lemma}
\label{lemma-coherent-closed}
Let $B$ be an algebraic space. Let $\pi : X \to Y$ be a closed immersion
of algebraic spaces which are separated and of finite presentation over $B$.
Then the morphism $\Cohstack_{X/B} \to \Cohstack_{Y/B}$ of
Lemma \ref{lemma-coherent-functorial} is a closed immersion.
\end{lemma}

\begin{proof}
Let $\mathcal{I} \subset \mathcal{O}_Y$ be the sheaf of ideals cutting
out $X$ as a closed subspace of $Y$. Recall that $\pi_*$ induces
an equivalence between the category of quasi-coherent $\mathcal{O}_X$-modules
and the category of quasi-coherent $\mathcal{O}_Y$-modules annihilated
by $\mathcal{I}$, see Morphisms of Spaces, Lemma
\ref{spaces-morphisms-lemma-i-star-equivalence}.
The same, mutatis mutandis, is true after base by $T \to B$ with
$\mathcal{I}$ replaced by the ideal sheaf
$\mathcal{I}_T = \Im((Y_T \to Y)^*\mathcal{I} \to \mathcal{O}_{Y_T})$.
Analyzing the proof of Lemma \ref{lemma-coherent-functorial}
we find that the essential image of
$\Cohstack_{X/B} \to \Cohstack_{Y/B}$
is exactly the objects $\xi = (T \to B, \mathcal{F})$
where $\mathcal{F}$ is annihilated by $\mathcal{I}_T$.
In other words, $\xi$ is in the essential image if and only if
the multiplication map
$$
\mathcal{F} \otimes_{\mathcal{O}_{Y_T}} (Y_T \to Y)^*\mathcal{I}
\longrightarrow
\mathcal{F}
$$
is zero and similarly after any further base change $T' \to T$.
Note that
$$
(Y_{T'} \to Y_T)^*(
\mathcal{F} \otimes_{\mathcal{O}_{Y_T}} (Y_T \to Y)^*\mathcal{I}) =
(Y_{T'} \to Y_T)^*\mathcal{F} \otimes_{\mathcal{O}_{Y_{T'}}}
(Y_{T'} \to Y)^*\mathcal{I})
$$
Hence the vanishing of the multiplication map on $T'$
is representable by a closed subspace of $T$ by
Flatness on Spaces, Lemma \ref{spaces-flat-lemma-F-zero-closed-proper}.
\end{proof}

\end{document}
